\documentclass[letterpaper,12pt,twoside,]{pinp}

%% Some pieces required from the pandoc template
\providecommand{\tightlist}{%
  \setlength{\itemsep}{0pt}\setlength{\parskip}{0pt}}

% Use the lineno option to display guide line numbers if required.
% Note that the use of elements such as single-column equations
% may affect the guide line number alignment.

\usepackage[T1]{fontenc}
\usepackage[utf8]{inputenc}

% pinp change: the geometry package layout settings need to be set here, not in pinp.cls
\geometry{layoutsize={0.95588\paperwidth,0.98864\paperheight},%
  layouthoffset=0.02206\paperwidth, layoutvoffset=0.00568\paperheight}

\definecolor{pinpblue}{HTML}{185FAF}  % imagecolorpicker on blue for new R logo
\definecolor{pnasbluetext}{RGB}{101,0,0} %


\usepackage{booktabs}
\usepackage{longtable}
\usepackage{array}
\usepackage{multirow}
\usepackage{wrapfig}
\usepackage{float}
\usepackage{colortbl}
\usepackage{pdflscape}
\usepackage{tabu}
\usepackage{threeparttable}
\usepackage{threeparttablex}
\usepackage[normalem]{ulem}
\usepackage{makecell}

\title{Assignment 3 - One sample rates and paramter contrasts. Due November 15,
11:59pm 2020}

\author[a]{EPIB607 - Inferential Statistics}

  \affil[a]{Fall 2020, McGill University}

\setcounter{secnumdepth}{5}

% Please give the surname of the lead author for the running footer
\leadauthor{Bhatnagar}

% Keywords are not mandatory, but authors are strongly encouraged to provide them. If provided, please include two to five keywords, separated by the pipe symbol, e.g:
 \keywords{  Rates |  Parameter contrasts |  Regression  }  

\begin{abstract}
In this assignment you will practice one sample rates and parameter
contrasts with linear regression. State all assumptions. Provide
confidence intervals with appropriate units. Answers should be given in
full sentences (DO NOT just provide the number). All graphs and
calculations are to be completed in an R Markdown document using the
provided template (from A1). You are free to choose any function from
any package to complete the assignment. Concise answers will be
rewarded. Be brief and to the point. Please submit the compiled pdf
report to Crowdmark by November 15, 2020, 11:59pm. You need to save your
answers to each question in separate pdf files. You also need to upload
your code separately to Q5. See \url{https://crowdmark.com/help/} for
details.
\end{abstract}

\dates{This version was compiled on \today} 

% initially we use doi so keep for backwards compatibility
% new name is doi_footer


\begin{document}

% Optional adjustment to line up main text (after abstract) of first page with line numbers, when using both lineno and twocolumn options.
% You should only change this length when you've finalised the article contents.
\verticaladjustment{-2pt}

\maketitle
\thispagestyle{firststyle}
\ifthenelse{\boolean{shortarticle}}{\ifthenelse{\boolean{singlecolumn}}{\abscontentformatted}{\abscontent}}{}

% If your first paragraph (i.e. with the \dropcap) contains a list environment (quote, quotation, theorem, definition, enumerate, itemize...), the line after the list may have some extra indentation. If this is the case, add \parshape=0 to the end of the list environment.


\hypertarget{points-2-each-seismicity-before-and-after-hydraulic-fracturing-in-the-horn-river-basin-northeast-british-columbia}{%
\section{(12 points, 2 each) Seismicity before and after hydraulic
fracturing in the Horn River Basin, northeast British
Columbia}\label{points-2-each-seismicity-before-and-after-hydraulic-fracturing-in-the-horn-river-basin-northeast-british-columbia}}

Consult the article \emph{Investigation of regional seismicity before
and after hydraulic fracturing in the Horn River Basin, northeast
British Columbia} by Farahbod et al.~(2014).

\begin{enumerate}
\def\labelenumi{\alph{enumi})}
\item
  Using just the year 2006 data in the first four columns of Table 2,
  calculate separate event rates for the periods when hydraulic
  fracturing (i) was and (ii) was not in operation. Express each rate as
  the number of events per year, and accompany it with a 95\% CI.
\item
  Repeat the calculations for the individual years 2007 to 2011.
\item
  Repeat the calculations for the entire period 2006 to 2011.
\item
  Store your results in a tidy \texttt{data.frame}.
\item
  Are you comfortable using Poisson-based CIs for the entire period 2006
  to 2011? i.e.~does it look like the rate in non-HF days is homogeneous
  over the 6 years, i.e., do the 6 CIs largely overlap each other?
\item
  Does it look like the rate in HF days is homogeneous over the 6 years,
  i.e., do the 6 CIs largely overlap each other?
\end{enumerate}

\hypertarget{points-simulation-study-for-confidence-intervals-of-counts}{%
\section{(28 points) Simulation study for confidence intervals of
counts}\label{points-simulation-study-for-confidence-intervals-of-counts}}

The goal of this simulation study is to estimate the coverage
probabilties of three different confidence intervals for a Poisson
count. In class we saw at least three methods to calculate the
confidence interval for a count: 1) normal approximation, 2) the exact
method and 3) using the Poisson quantile \texttt{stats::qpois} function.

\begin{enumerate}
\def\labelenumi{\alph{enumi})}
\item
  (12 points) Simulate 1000 trials from a Poisson(\(\mu\) = 5)
  distribution using the \texttt{stats::rpois} function. For each of
  these trials, calculate the 95\% confidence interval for the count
  using each of the three methods mentioned above. For each of the three
  methods, calculate the coverage probability, i.e., the percentage of
  intervals which contain the true mean \(\mu\). Describe your findings
  and comment on the differences between methods in terms of coverage
  probability.
\item
  (16 points) Repeat the simulation study in a), but this time, with
  increasing values of \(\mu\). Visualize the coverage probabilities as
  a function of \(\mu\) for each of the three methods. Describe your
  findings and comment on the differences between methods in terms of
  coverage probabilities as a function of \(\mu\), i.e., what happens to
  the coverage probabilities when the expected number of counts
  increases? Explain.
\end{enumerate}

\hypertarget{points-5-each-weather-data-analysis}{%
\section{(30 points, 5 each) Weather Data
Analysis}\label{points-5-each-weather-data-analysis}}

Environment and Climate Change Canada (ECCC), formerly Environment
Canada, is a federal department with the stated role of protecting the
environment, conserving national natural heritage, and also providing
weather and meteorological information. According to ECCC \emph{warming
over the 20th century is indisputable and largely due to human
activities} adding \emph{Canada's rate of warming is about twice the
global rate: a 2°C increase globally means a 3 to 4ºC increase for
Canada}. Berkeley Earth has reported that
\href{http://berkeleyearth.org/wp-content/uploads/2016/01/2015-Hottest-Year-BE-Press-Release-v1.0.pdf}{2015
was unambiguously the warmest year on record across the world}, with the
Earth's temperature more than 1.0 C (1.8 F) above the 1850-1900
average.\\
The data set \texttt{qc\_weather.csv} contains daily temperature
readings at
\href{https://www.wunderground.com/weather/ca/danville}{Danville weather
station in Quebec} starting from June 6, 1880 (some dates are missing)
that were scraped from the
\href{https://climate.weather.gc.ca/historical_data/search_historic_data_e.html}{Environment
and Climate Change Canada (ECCC) website} using the
\href{https://docs.ropensci.org/weathercan/index.html}{weathercan R
package}. The daily measurements provided for you are:

\begin{itemize}
\item
  \texttt{year}: the year of the measurement.
\item
  \texttt{month}: the month of the measurement; 1 = January, 2 =
  February,\ldots{} , 12 = December.
\item
  \texttt{day}: the day within the month (1-31).
\item
  \texttt{max\_temp}: the maximum recorded temperature that day.
\item
  \texttt{mean\_temp}: the average of recorded temperatures that day.
\item
  \texttt{min\_temp}: the minimum recorded temperature that day.
\item
  \texttt{dayfromstart}: a number from 1 to 13501 that represents the
  number of days since June 5. So 1 = June 6, 1880.
\end{itemize}

\begin{enumerate}
\def\labelenumi{\alph{enumi})}
\tightlist
\item
  Create a new data set which contains all the rows that correspond to
  February in any year. Compare the \texttt{max\_temp} in the February
  data between (1) those February days in 2015 and (2) February days in
  all other years using an appropriate figure and summary statistics.
  Describe your findings in 3 or fewer sentences.
\end{enumerate}

\begin{enumerate}
\def\labelenumi{\alph{enumi})}
\setcounter{enumi}{1}
\item
  Were those February days in 2015 warmer than February days in all
  other years ? Use an appropriate regression model to answer this
  question. Write the regression equation in terms of population
  parameters and define all parameters in your model. What is the
  parameter of interest?
\item
  Estimate the regression parameters using the data provided. Report the
  estimated coefficient for the parameter of interest and interpret it
  in the context of the problem.
\item
  Report the p-value for the parameter of interest. Based on the
  p-value, is there sufficient evidence at \(\alpha = 0.05\) to support
  the claim that 2015 was the warmest year on record? What assumptions
  were used in calculating the p-value.
\item
  Provide a potential explanation for why your analysis agrees or does
  not agree with the claim made by Berkeley Earth.
\item
  Use a permutation test to calculate the p-value for the parameter of
  interest and compare it with the one obtained in part d). Briefly
  discuss this comparison.
\end{enumerate}

\hypertarget{points-6-each-physical-activity-in-nhanes}{%
\section{(30 points, 6 each) Physical activity in
NHANES}\label{points-6-each-physical-activity-in-nhanes}}

This problem uses data from the
\href{https://cran.r-project.org/web/packages/NHANES/NHANES.pdf}{National
Health and Nutrition Examination Survey (NHANES)}, a survey conducted
annually by the US Centers for Disease Control (CDC). The dataset is
available from the \texttt{NHANES} package.

Regular physical activity is important for maintaining a healthy weight,
boosting mood, and reducing risk for diabetes, heart attack, and stroke.
In this problem, you will be exploring the relationship between weight
(\texttt{Weight}) and physical activity (\texttt{PhysActive}) using the
NHANES data. Weight is measured in kilograms. The variable
\texttt{PhysActive} is coded \texttt{Yes} if the participant does
moderate or vigorous-intensity sports, fitness, or recreational
activities, and \texttt{No} if otherwise. The objective is to compare
weight between physically active and those who are not.

\begin{enumerate}
\def\labelenumi{\alph{enumi})}
\item
  Explore the data.

  \begin{enumerate}
  \def\labelenumii{\roman{enumii}.}
  \item
    Identify how many individuals are physically active.
  \item
    Create a plot that shows the association between weight and physical
    activity. Describe what you see.
  \end{enumerate}
\item
  Provide an appropriate regression model for the stated objective and
  state the parameter of interest. Give the regression equation in terms
  of population parameters and be sure to define each of the parameters
  in your model.
\item
  Fit a linear regression model to estimate the regression parameters.
  Report the estimated coefficients from the model and interpret each of
  them in the context of the problem.
\item
  Report a 95\% confidence interval for the parameter of interest and
  interpret the interval in the context of the problem. Based on the
  interval, is there sufficient evidence at \(\alpha = 0.05\) to reject
  the null hypothesis of no association between weight and physical
  activity? State the assumptions used for calculating the 95\%
  confidence interval.
\item
  Provide a 95\% bootstrap confidence interval for the parameter of
  interest and compare it to the one in part d). Briefly discuss the
  comparison.
\end{enumerate}

%\showmatmethods


\bibliography{pinp}
\bibliographystyle{jss}



\end{document}

